% Imported from topics/video-generation/seedance-2.5/reference-binding-spec.md
\chapter{主体槽、参考冲突与细节预算}\label{seedance-m22}
\moduleinfo{M22}{候选设计与工程选型；组合效果待验证}{\hyperref[src-alchemint]{\texttt{alchemint}}、\hyperref[src-phantom]{\texttt{phantom}}、\hyperref[src-ltx-2]{\texttt{ltx-2}}、\hyperref[src-ltx-original-code]{\texttt{ltx-original-code}}、\hyperref[src-vace]{\texttt{vace}}}

参考的内容、归属与时间已经定义，接着把素材身份、池化预算和读取规则落实到同一接口。

再看一个完整参考输入：角色 A 有六张脸和服装图，角色 B 有两张，商品图约束外观，动作来自一段视频，声音又来自两段语音。文件都能编码，并不说明生成 A 时会读到 A 的脸，或动作视频中的衣服不会一并复制。还需要把这些关系表示出来。

所以本节继续固定素材、主体、属性和时间之间的对应，再讨论冲突与预算。以下仍是联合音视频主干上的候选规则，宽度和压缩容量尚未经模型验证。

\section{主体槽是一次任务中的关系}\label{seedance-m22-ux4e3bux4f53ux69fdux662fux4e00ux6b21ux4efbux52a1ux4e2dux7684ux5173ux7cfb}

素材身份用内容哈希和 \codepath{asset_id}；人物或商品用稳定业务ID；模型则使用场景局部主体槽，例如A、B、product。业务ID用于检索，局部槽用于表示这一次请求中的对应关系。

局部槽号的作用是本次请求中的对应，并不是永久身份。训练时可同时置换 A/B 槽，以及所有相关文本、图像、声音与区域标签，迫使模型读关系而不是记住第 0 槽总是主角。人物实际长相仍由素材给出。

\begin{verbatim}
{
  "asset_id": "a_coat_front",
  "subject_slot": "A",
  "roles": ["clothing"],
  "shot_ids": ["s2"],
  "active_intervals_seconds": [[8, 12], [15, 18]],
  "priority": 2,
  "confidence": 0.95
}
\end{verbatim}

同一图片需要同时提供脸和服装时，保留一个资产、建立两个属性视图；不复制成两个互相无关的主体。音频另保留speaker-to-subject映射，群体环境声没有人物槽。

\section{实体词与视觉特征怎样共享绑定}\label{seedance-m22-ux5b9eux4f53ux8bcdux4e0eux89c6ux89c9ux7279ux5f81ux600eux6837ux5171ux4eabux7ed1ux5b9a}

实体文本先从冻结文本模型的原始 token 序列中定位描述范围，汇聚特征，再投影成 512 维记忆。与它对应的图像摘要使用相同主体槽和 role 表示，于是文字与视觉条件有了共同指向。

字符偏移还不能直接当作连接器后的 token 位置：连接器可能改变序列并添加 register。首版因此单独保存实体条件记忆，原 caption 仍沿原生文本路径，同时记录实体跨度和 tokenizer 版本，避免索引看似对齐，读到的却是另一段内容。

每个图像资产独立VAE编码，空间坐标与局部crop信息进入池化器；每段参考视频保留自己的时间轴，不与其它素材在像素时间维拼接。音色素材由音频codec独立编码，文字转录只用于区分声音内容与声音身份。

\section{先分配预算，再进行汇聚}\label{seedance-m22-ux5148ux5206ux914dux9884ux7b97ux518dux8fdbux884cux6c47ux805a}

候选共享记忆宽度512。视觉总预算2048，声音摘要预算512，null token单独增加一个位置；这些预算表示可读取记忆长度，不包含原始素材编码开销。

{\def\LTcaptype{table} % do not increment counter
\small\setlength{\tabcolsep}{3pt}\renewcommand{\arraystretch}{1.18}
\begin{longtable}[]{@{}
  >{\raggedright\arraybackslash}p{(\linewidth - 4\tabcolsep) * \real{0.1900}}
  >{\raggedleft\arraybackslash}p{(\linewidth - 4\tabcolsep) * \real{0.3900}}
  >{\raggedright\arraybackslash}p{(\linewidth - 4\tabcolsep) * \real{0.4200}}@{}}
\toprule\noalign{}
\begin{minipage}[b]{\linewidth}\raggedright
角色
\end{minipage} & \begin{minipage}[b]{\linewidth}\raggedleft
默认摘要数
\end{minipage} & \begin{minipage}[b]{\linewidth}\raggedright
额外保留什么
\end{minipage} \\
\midrule\noalign{}
\endhead
\bottomrule\noalign{}
\endlastfoot
人物/商品图片 & 32/图 & 主体crop、视角和关键细节crop \\
外观参考视频 & 64/段 & 分段时间、视角与可见性 \\
音色/声音风格 & 32/段 & speaker、干净语音段与文字内容标签 \\
仅动作/相机视频 & 不占身份摘要 & pose/flow/相机走结构路径，单独计费 \\
必须逐字/逐拍保留的音轨 & 不压成声音摘要 & 完整原生时序音频条件 \\
\end{longtable}
}

预算先满足本镜头必要主体，再按优先级和不确定性分配剩余量。不够时先去掉无关镜头、近重复视角和低质量素材，并登记裁减结果。否则唯一背面或唯一 speaker 参考被悄悄丢掉，后面再调强度也不能补回这份条件。

30图＋10外观视频按默认预算为1600视觉位置，10音频为320声音位置。两类都能放进预算，只证明输入可表示；没有证明小logo、多人身份和音色细节仍然充分。

池化前保留crop坐标、素材内时间和type信息。对细字、商标和精密结构，摘要容量不足时增加局部预算，或通过确定性图层交付；摘要token不能承诺恢复已经丢掉的像素信息。

\section{可读性由三层规则决定}\label{seedance-m22-ux53efux8bfbux6027ux7531ux4e09ux5c42ux89c4ux5219ux51b3ux5b9a}

对视频query \(i\) 与参考位置 \(j\)，先判定镜头和属性是否允许，再计算时间偏置：

\[
B_{ij}=\log\left(\varepsilon+(1-\varepsilon)g_j(\tau_i)\right)
+\log(\max(c_j,\varepsilon))+H_{ij}.
\]

\(c_j\in[0,1]\)是置信度；\(H_{ij}\)为兼容时0、不兼容时负无穷。时间函数 \(g_j\) 合并同主体的重叠区间，对多次出场分别计算软门，再限制到 \([0,1]\)。null位置始终有限且不施加置信度惩罚。

空间query没有已知主体槽，默认只做镜头级粗筛。用户布局、源视频跟踪或可靠预测区域存在时，可进一步限制空间归属；遮挡边界使用软区域，而不是强行把一个latent只分给一个人物。

多人混合音频的一个时间token也可能含两个speaker。声音支路按明确对白时间和说话人映射路由，重叠讲话使用集合或软权重；没有说话人时间信息时保留不确定性，不能由文件顺序猜声音归属。

\section{冲突如何进入任务编译}\label{seedance-m22-ux51b2ux7a81ux5982ux4f55ux8fdbux5165ux4efbux52a1ux7f16ux8bd1}

相似图像未必是冲突：同一人的两个视角可以互补；同一时刻又要求同一人穿两套互斥服装，才需要裁决。判断必须结合属性、时间与优先级，对全部图像简单平均没有表达这种区别。

{\def\LTcaptype{table} % do not increment counter
\small\setlength{\tabcolsep}{3pt}\renewcommand{\arraystretch}{1.18}
\begin{longtable}[]{@{}
  >{\raggedright\arraybackslash}p{(\linewidth - 2\tabcolsep) * \real{0.2200}}
  >{\raggedright\arraybackslash}p{(\linewidth - 2\tabcolsep) * \real{0.7800}}@{}}
\toprule\noalign{}
\begin{minipage}[b]{\linewidth}\raggedright
情况
\end{minipage} & \begin{minipage}[b]{\linewidth}\raggedright
首版处理
\end{minipage} \\
\midrule\noalign{}
\endhead
\bottomrule\noalign{}
\endlastfoot
同主体多视角 & 共用主体槽，保留视角标签 \\
同主体前后换衣 & 为服装属性分别设置区间，身份参考持续有效 \\
同一时段两套互斥服装 & 采用显式优先级；无优先级则返回规格冲突 \\
动作视频人物与身份参考不同 & 禁用其外观记忆，只传结构条件 \\
声音参考包含音乐/他人讲话 & 分段或分离后用；无法分离则降低置信度或回到后期音轨方案 \\
同一角色离场后返回 & 保留身份记忆，出场监督记录多个区间 \\
\end{longtable}
}

置信度表示我们是否确定这条信息，并不改变用户是否要求它成立。低置信但必需的条件，应回到输入修正，而不能被自动降权后当作已经满足。

\section{训练要教会什么}\label{seedance-m22-ux8badux7ec3ux8981ux6559ux4f1aux4ec0ux4e48}

参考正确但背景不同、参考服装不同、相似人物同时出现、人物暂时遮挡、角色真正离场、声色与台词不同，是不同的数据切片。只收集``人物一直在画面中央''的正例，不足以学习绑定。

保持目标不变时，可以同时置换槽号或素材顺序，也可以增加无关素材、去掉重复视角，检查模型是否仍遵守同一关系。但把 A 的参考换成 B，目标仍写 A，就已经改变条件与目标的匹配，不能继续当成同一条正确生成监督。

错误参考可用来训练冲突检测、绑定分类或偏好比较；生成分支则只作训练约定内的条件 dropout。这样模型学到的“忽略条件”才有明确原因，不会被矛盾样本反复强化。

\section{验收使用成组任务}\label{seedance-m22-ux9a8cux6536ux4f7fux7528ux6210ux7ec4ux4efbux52a1}

至少保留单主体新背景、两主体相似外观、多人重叠、换衣、间断出场、动作换身份、同音色新台词和文件数上限八类。每类分别看身份、动作、误出场、边界偏差、背景泄漏和声音归属，不能只报告总体美感。

几个对照可以直接检验关系是否真的被使用：文本不动只换参考，参考不动只换时间，所有槽同步置换，或者只打乱文件顺序。正确行为应随语义绑定变化，而不应依赖某个文件刚好排在第几位。

时间偏置的极简示例见 \hyperref[code-reference-timing]{reference\_timing.py}，候选读取见 \hyperref[code-gated-condition]{gated\_condition.py}。真实主体参考结构可读 \href{https://github.com/Phantom-video/Phantom}{Phantom}，它提供的实现和这里的摘要路由、音频绑定不是同一套完整工程。
